\chapter{Explainability}

\section{Introduction}

\subsection{Why Explainability?}

In recent years, artificial intelligence has advanced at a dizzying pace, with models becoming increasingly accurate, complex, and difficult to interpret. Algorithms such as deep neural networks, ensembles, and generative models offer outstanding performance, but their opacity poses a critical problem: how can we trust a system we cannot understand?

\textbf{Explainability} seeks to address this need. Its goal is to make AI models understandable to humans, whether end users, developers, or business leaders. This not only improves confidence in the results but also allows for the detection of errors, biases, and vulnerabilities that might otherwise go unnoticed.

The motivation behind explainability goes beyond technical curiosity. In sectors such as medicine, finance, and law, an automated decision without clear justification can have serious ethical, legal, and social consequences. Furthermore, with regulations such as the General Data Protection Regulation (GDPR) in Europe, the ability to explain automated decisions is no longer just desirable, but mandatory.

This chapter explores why explainability is key to developing responsible and safe AI models, what techniques are currently available, and how to choose the best strategy based on the context of the problem.

\subsection{Interpretability vs Explainability}

Although the terms interpretability and explainability are often used interchangeably, in the context of artificial intelligence they represent distinct and complementary concepts.

\textbf{Interpretability} refers to the ability to directly understand how a model works. That is, a person can inspect the model and reason about its behavior without the need for external tools. For example, a small decision tree or a linear regression are interpretable because their rules and coefficients can be followed step by step. This type of models are called \textbf{white-box} models.

\textbf{Explainability}, on the other hand, focuses on generating explanations about model predictions, even if the models are not directly interpretable. For example, a deep neural network may be inexplicable on its own, but using some techniques, we can construct explanations that help us understand why the model made a certain decision. This type of models are called \textbf{black-box} models. 

In this chapter, we will study explainability techniques for black-box models, as they are the most used in computer vision.

\section{Visualization of Internal Features}

\subsection{Maximally Activating Patches}

A common technique for interpreting convolutional networks is to analyze \textbf{what types of patterns cause certain layers or filters to activate more strongly}. This is known as searching for maximally activating patches, that is, the image fragments that provoke the greatest activation in a specific layer/channel of a CNN layer.

Note that when we talk about the activations of a channel we refer to a channel of a specific layer and when we talk about the activations of a layer we refer to all the channels of that layer (we can, for example, take the mean or maximum along the channels dimension to obtain a 2D activation).

\begin{figure}[H]
\centering
\includegraphics[scale=0.38]{images/chapter_5/max_act_pat}
\caption{Maximally activating patches of different layers in a CNN. Layers 14 and 17 visualize the eyes of the cat. Layers 19 and 21 visualize the paws, and layer 24, 26, and 28 visualize the face.}
\end{figure}

The idea is simple: if a layer/channel responds strongly to a specific region of an image, that region contains information that the network considers relevant. This technique allows to visualize and understand what the different layers/channels represent internally, helping to validate that the model is learning features consistent with the problem, or to discover if it is focusing on irrelevant or biased details.

The algorithm is the following:
\begin{enumerate}
\item Run many images through the network and record the activations of chosen channel/layer.
\item Reshape the activations to the shape of the original image.
\item Visualize image patches that correspond to maximal activations.
\end{enumerate}

Maximally activating patches do not explain all of a CNN's behavior, but they offer a valuable window into the network's ``inside'' and are especially useful for auditing complex models and debugging their behavior.

\subsection{Gradient Ascent}

A powerful technique for understanding \textbf{what activates a specific layer or filter} in a CNN is gradient ascent. Unlike maximally activating patches, which look for real-world examples that trigger high activations, gradient ascent generates synthetic images that directly maximize the activation of a neuron, channel, or class.

The procedure is as follows: starting with an initial image (it can be noise or a base image), its pixels are iteratively adjusted using the gradient of the target activation with respect to the input. In other words, it applies an inverted version of traditional training: instead of modifying the model's weights to minimize a loss, the image pixels are modified to maximize the desired activation.

Mathematically, this can be seen as:
$$\text{argmax}_x \{\text{agg}\left(f(x)\right)-\lambda ||x||_2\}=\text{argmax}_x\{\mathcal{F}(x)\},$$
where $x$ is the image we want to find, $f(x)$ the activation of the layer/channel and  $\lambda$ a parameter to penalize the $L_2$ norm of the generated image (regularization term). The term agg just means an aggregation function:
\begin{itemize}
\item If $f$ refers to the final output for a class, $\text{agg}\left(f(x)\right)=f(x)$, as $f(x)$ is a scalar and we can not apply an aggregation function.
\item If $f$ refers to the activation of a channel of a layer, agg can be the mean, maximum or the norm along height and width of the activation.
\item If $f$ refers to the activation of a layer, agg can be the mean, maximum or the norm along channels, height and width of the activations of the layer.
\item In some cases, it is also interesting to take the aggregation function of the absolute values. This is because large, negative terms can also be important.
\end{itemize}
To solve this problem, we can repeat for many iterations
$$x_{t+1} = x_t + \alpha \Delta \mathcal{F}(x_t),$$
where $x_t$ is the image at iteration $t$, $\Delta \mathcal{F}(x_t)$ its gradient with respect to the function and $\alpha$ the learning rate.

The result of this process is a synthetic image that represents the ``ideal stimulus'' for a specific unit of the model. For example, we can generate an image that maximizes the activation of filter 23 in the third convolutional layer, or the output corresponding to the ``cat'' class in the final layer.

These visualizations allow us to understand what kind of abstract patterns the model has learned in different parts of the network, although the result can often be difficult to interpret directly (due to its artificial nature). To improve readability, techniques such as regularization or clipping are often applied to stabilize the process and generate more natural images.

\begin{figure}[H]
\centering
\includegraphics[scale=0.40]{images/chapter_5/gascent}
\caption{Images generated for maximizing the final score of different classes with gradient ascent.}
\end{figure}

\subsection{Feature Inversion}

Feature inversion is a technique that seeks to answer the following question: \textbf{what information is the network actually storing in its internal representations?} To do this, it attempts to reconstruct an input image from its activations in a specific intermediate layer/channel of the network.

The procedure is conceptually simple: an image is taken, its internal representation is calculated (e.g., the output of convolutional layer 4), and then an attempt is made to find a new image that, when passed through the network, produces activations as similar as possible to the original ones. In other words, an attempt is made to reverse the feature extraction process.

\begin{figure}[H]
\centering
\includegraphics[scale=0.52]{images/chapter_5/inverse}
\caption{Original images (first column) and reconstructed images using feature inversion (second column).}
\end{figure}

This is formulated as an optimization problem: starting with an initial image (noise or a base image), its pixels are adjusted to minimize the difference between the activations generated by that image and the target activations. Mathematically, we want to minimize:
$$\mathcal{F}(x) = ||f(x)-f(x_0)||_2^2 + \lambda ||x||_2,$$
where $x_0$ is the original image, $f(x)$ the activation of the layer/channel and we use again $\lambda$ as a regularization term. A similar iterative process as in gradient ascent can be followed to solve the problem.

The result of this process is a new image that reflects what the network ``remembers'' about the input at a given level of abstraction. At lower layers, reconstructions are typically very similar to the original image (with clear edges and textures), while at deeper layers, the reconstructed image becomes more fuzzy and semantic, losing specific details but preserving the overall structure.

This technique reveals how the network discards unnecessary information as it goes deeper, retaining only what is relevant to the task (e.g., general shape, position, key texture). Feature inversion is useful for auditing models, detecting whether they are memorizing irrelevant details, and understanding the degree of compression or generalization occurring at each layer of the CNN.

\subsection{DeepDream}

DeepDream is a technique developed by Google engineers that explores \textbf{what a neural network ``sees'' when asked to boost what it already recognizes}. Although it emerged as an artistic tool, it has also proven useful for visualizing what internal patterns activate certain neurons or layers.

The idea behind DeepDream is similar to gradient ascent, but with a key difference: instead of maximizing a specific class or activation from scratch, it starts with a real image and amplifies the activations already present. It is like telling the network: ``Take what you already recognize in this image and make it stronger''.

The process to solve the problem is exactly the same as in gradient ascent, but now the initial image is the original image instead of noise or a base image.

As the iterative process is repeated, the image begins to transform and display patterns that the network ``dreams'' or amplifies: eyes, feathers, buildings, animals, etc., depending on what the network has learned during training. Hence the name: DeepDream, as if the network were imagining or hallucinating what it sees.

\begin{figure}[h]
    \centering
    \begin{minipage}{0.48\textwidth}
        \centering
        \includegraphics[width=\textwidth]{images/chapter_5/dfirst}
    \end{minipage}
    \hfill
    \begin{minipage}{0.48\textwidth}
        \centering
        \includegraphics[width=\textwidth]{images/chapter_5/dlast}
    \end{minipage}
    \caption{DeepDream applied to one of the first layers (left) and one of the last layers (right). Note that in the first image more textures are remarked, while in the second image more high level objects are remarked.}
\end{figure}

In terms of interpretability, this technique does not offer precise or quantitative explanations, but it is useful for gaining a visual intuition of the latent concepts the network has internalized. It is also an effective way to illustrate how a CNN represents knowledge in multiple hierarchical layers.

\section{Perturbation-Based Methods}

Perturbation-based methods are explainability techniques based on a simple but powerful idea: \textbf{hide parts of the input and observe how the output of the model changes}. If hiding a specific region significantly affects the prediction, then that region is important for the model's decision.

In the case of images, this is typically implemented by covering blocks (patches) of pixels with a neutral value (such as gray or black), one by one, and evaluating how the probability of the target class changes. The result is a heat map showing which regions of the image were most relevant to the prediction.

This technique is completely independent of the model architecture and does not require access to gradients, which is why it is considered a \textbf{model-agnostic} method. This means we can apply these techniques to any model: linear regression, random forest, neural networks\ldots It is also intuitive and easy to visualize, although \textbf{computationally expensive} if many occlusions are applied.

One common problem is the \textbf{size of the patch}, also known as superpixel, we hide. Here are some considerations:
\begin{itemize}
\item The bigger the patch, the less computationally expensive the method will be, as there will be less forward passes.
\item The smaller the patch, the the more accurate contribution of the patch, but the method will be more computationally expensive, as there will be more forward passes. Also, if the patch is too small, it may not capture the local context of the patch.
\end{itemize}
There is not a optimal solution for this problem. However, we have to take into account the pros and cons of each patch size.

\subsection{SHAP}

\subsubsection{General Idea}

\textbf{SHAP} (SHapley Additive ExPlanations) is a \textbf{game-theoretic} explainability method that assigns each input feature a quantitative contribution to the model's prediction. Although originally designed for tabular data, it can also be applied to images, with specific adaptations. In the case of images, each ``feature'' can be an individual pixel or a group of pixels (e.g., a superpixel or a patch). The goal of SHAP is to estimate how much each image region contributes to the model output for a specific class.

\begin{figure}[H]
\centering
\includegraphics[scale=0.75]{images/chapter_5/shap}
\caption{SHAP applied for different classes in two images.}
\label{shap}
\end{figure}

The algorithm is the following:
\begin{enumerate}
\item The image is divided into regions, usually square groups of pixels.
\item All possible versions of the image are created where some regions are present and others are hidden (as if those features did not exist). They are replaced usually with zeros.
\item For each combination of visible/invisible regions, the model output is calculated.
\item Calculation of SHAP values: Using these outputs, the individual contribution of each region is estimated, based on the weighted average of its impact over all possible combinations (Shapley values).
\end{enumerate}


It produces heat map where each region has a SHAP value indicating whether its presence increases the probability of the target class (positive contribution, usually in red) or decreases the probability (negative contribution usually in blue). This map allows you to see which parts of the image support or contradict the model's prediction. An example is shown in Figure \ref{shap}.

\subsubsection{Calculation of Shapley Values}

Let's assume we divide the image in $M$ regions and we want to calculate the contribution of each region to the prediction for a specific class of the model $f$. Let's also denote $N=\{1, \ldots, M\}$. For each subset of regions $S\subseteq N-\{i\}, i = 1, \ldots, M$, the model prediction is evaluated when only the regions of $S$ are visible, and it is compared to the prediction when region $i$ is also included. This part corresponds to the absolute value (right term) of the formula.

The general formula for the SHAP value for region $i$ is:
$$\phi_i = \sum_{S\subseteq N-\{i\}} \dfrac{|S|! \left(M-|S|-1\right)!}{M!}\left|f\left(S\cup \{i\}\right) - f\left(S\right)\right|.$$
The combinatorial coefficient (the fraction, first term of the formula) is the Shapley weight, which guarantees fairness in the assignment (each order of appearance has equal probability). SHAP values are based on the idea that each possible order of appearance of regions is equally likely. Therefore, to calculate the average contribution of a region $i$, we must consider all possible positions in which it could be added and average their marginal contributions fairly.

\subsubsection{Example}

Suppose the regions are $A$, $B$ and $C$. Let's calculate the Shapley weight for $A$. We have that:
\begin{itemize}
\item $S\subseteq N-\{A\} \equiv S_0=\emptyset, S_1=\{B\}, S_2=\{C\}, S_3=\{B, C\}$. Note that $\{C, B\}$ is not included because the order does not matter.
\item The weight for $S_0$ is $w(S_0)=0!(3-0-1)!/3!=1/3$.
\item Using the same formula, $w(S_1)=1/6$, $w(S_2)=1/6$ and $w(S_3)=1/3$.
\end{itemize}
An intuitive way of understanding the weight is writing all possible permutations:
\begin{itemize}
\item ABC
\item ACB
\item BAC
\item BCA
\item CAB
\item CBA
\end{itemize}
We can deduce:
\begin{itemize}
\item The probability that the first is $A$ is $2/6=1/3$, so $w(S_0)=1/3$.
\item The probability that the second is $A$ and the first is $B$ is $1/6$, so $w(S_1)=1/6$.
\item The probability that the second is $A$ and the first is $C$ is $1/6$, so $w(S_2)=1/6$.
\item The probability that the third is $A$ and the first and second are, $B$ and $C$ or $C$ and $B$, is $2/6=1/3$, so $w(S_3)=1/3$. 
\end{itemize}

\subsubsection{What is it Done in Practice?}

In practice, generating so many combinations may be prohibitive. Therefore, in some cases, the following approach is followed:
\begin{enumerate}
\item A new dataset of images is generated using a sample of all possible subsets $S$ and the prediction of the model is calculated for each subset. The bigger the sample, the more accurate the estimated Shapley values, but the more computationally expensive the method will be.
\item Each sample is weighted with the Kernel Shap:
$$\pi(S)=\dfrac{M-1}{\binom{M}{S}|S|(M-|S|)}.$$
This weight favors subsets with few or almost all active superpixels, which contribute more marginal information. This step is not compulsory, but it usually helps.
\item A linear regression is trained. The input is a one-hot vector of length $M$ with a 1 in position $i$ if the superpixel $i$ is present in $S$ and 0 otherwise. The target is the prediction of the original model. Note that the linear regression will have the form $\phi_0+\sum_{i=1}^M\phi_ix_i$.
\item The Shapley value of each region $i$ is $\phi_i$.
\end{enumerate}


\subsection{Occlusion}

\textbf{Occlusion} is a very intuitive technique that masks parts of the image systematically, calculates the difference with the original prediction and finally takes the mean along all channels.

\begin{figure}[H]
\centering
\includegraphics[scale=0.18]{images/chapter_5/occlusion}
\caption{Masked pixels for each iteration in occlusion.}
\label{occlusion}
\end{figure}

The algorithm is the following and a representation is show in Figure \ref{occlusion}:
\begin{enumerate}
\item Select a patch size, let's say $S$.
\item Initialize two tensors filled with zeros with the same shape as the original image $H\times W$. We will call them ``seen pixels'' and ``occlusion values''.
\item Mask pixels $[:]\times [0:S]\times [0:S]$ of the original image and calculate the new prediction of the model $f'(x)$.
\item Add 1 to positions $[0:S]\times [0, S]$ of ``seen pixels'' and $|f(x)-f'(x)|$ (difference between the original prediction and the new one) to the same positions of ``occlusion values''.
\item Repeat steps 3 and 4 for positions
$$[:]\times [0:S]\times[1:S+1], \ldots, [:]\times[0:S]\times[W-S:W],$$
$$[:]\times [1:S+1]\times[0:S], \ldots, [:]\times[1:S+1]\times[W-S:W],$$
$$\vdots$$
$$[:]\times [H-S:H]\times[0:S], \ldots, [:]\times [H-S:H]\times[W-S:W].$$
\item The output is the element-wise division of ``occlusion values'' by ``seen pixels''. This can be interpreted as the mean change per pixel.
\end{enumerate}
You can see that it is similar to a convolution. In fact, we can even modify the stride (it does not have necessarily to be 1) and some other parameters, but the original method is as described above. Finally, an example of an explanation obtained with Occlusion is shown in Figure \ref{Occlusion}. Note that the more granularity we want, the smaller the patch size will be, and more forward passes will have to be done. In some cases it can be prohibitive.

\begin{figure}[H]
    \centering
    \begin{minipage}{0.43\textwidth}
        \centering
        \includegraphics[width=\linewidth]{images/chapter_5/perro}
    \end{minipage}
    \hfill
    \begin{minipage}{0.47\textwidth}
        \centering
        \includegraphics[width=\linewidth]{images/chapter_5/Occlusion}
    \end{minipage}
\caption{Occlusion applied to a dog. The original image is on the left and the explanation map on the right.}
\label{Occlusion}
\end{figure}

\subsection{RISE}

\textbf{RISE} (Randomized Input Sampling for Explanation) is an explainability method whose key focus is using \textbf{random binary masks} and observing how they affect the model output.

The \textbf{algorithm} is the following  and a representation is shown in Figure \ref{rise}:
\begin{enumerate}
\item Initialize two tensors filled with zeros with the same shape as the original image $H\times W$. We will call them ``seen pixels'' and ``RISE values''.
\item Generate a random binary mask.
\item The original image is multiplied by the mask along all the channels, randomly hiding parts of the image.
\item Compute the output of the model for the new image $f'(x)$.
\item Add 1 to \underline{not} masked positions in ``seen pixels'' and $f'(x)$ to the same positions of ``RISE values''.
\item Repeat steps 2, 3 and 4 several times.
\item The output is the element-wise division of ``RISE values'' by ``seen pixels''. This can be interpreted as the mean contribution per pixel.
\end{enumerate}
Note that this method is very similar to occlusion, but here we do not use the original prediction because the values are interpreted as contributions of the pixels rather than changes caused by masked pixels. Therefore, instead of suming on masked values, we sum on not masked values.

For the \textbf{generation of the masks} many techniques can be applied, for example just creating a $H\times W$ tensor following a Bernoulli distribution with probability $p=0.5$. However, the following steps are usually followed to create $N$ masks:
\begin{enumerate}
\item Create a $S\times S$ mask following a Bernoulli$(p)$ distribution, in which each position is 1 with probability $p$ and 0 with probability $1-p$.
\item The mask is upsampled to the original image size using bilinear interpolation, which smooths the edges and makes the effect on the pixels smooth.
\item Apply random shifts (height and width offsets) to each mask to further diversify the masks and prevent them from all being rigidly aligned.
\end{enumerate}
Although the base is a per-pixel Bernoulli distribution in the initial mask, after upsampling and shifting, the result is a distribution of continuous and smooth masks, where each pixel in the final mask takes real values between 0 and 1. This allows for the calculation of more stable and smoother importance maps. Some common parameters are $N=4000$; $s=7$ or $s=14$; $p=0.5$; bilinear interpolation is used for upsampling, but any method can be used; and the shift range is uniform across the entire image, but it can also be changed.

\begin{figure}[H]
\centering
\includegraphics[scale=0.95]{images/chapter_5/rise}
\caption{RISE applied to a single image.}
\label{rise}
\end{figure}


\section{Gradient-Based Methods}

\textbf{Gradient-based} methods seek to explain a prediction by examining how the model's output varies when the input values are slightly perturbed. In other words, they measure the sensitivity of the output to each pixel in the image.

The central idea is simple: if a small change in a pixel produces a large change in the output (for example, in the probability of a class), that pixel is important for the decision.

They are \textbf{model-specific} methods, that is, they can only be applied to neural networks. However, they are much more \textbf{efficient} than occlusion methods, as usually only one forward pass and one backward pass is needed.

\subsection{Saliency Maps}

The \textbf{Saliency Map} was one of the first methods proposed to visualize which pixels in an image are most relevant to a neural network's decision. The idea is to calculate how sensitive the model's output is to each input pixel.

Given a classification model $f(x)$ and an input $x$ (e.g., an RGB image) with shape $3\times H\times W$, the saliency map for class $c$ is defined as the absolute value of the maximum value of the gradient of the output for class $c$ with respect to each pixel in the image along the channels dimension:
$$S_{i, j}(x)=\max_{k\in\{1, 2, 3\}}\left|\dfrac{\partial f_c(x)}{\partial x_{i, j, k}}\right|.$$
Note that $x_{i, j, k}$ is the pixel value at row $i$, column $j$ and channel $k$ and $S(x)\in\mathbb{R}^{H\times W}$.

This produces a map the same size as the image (with just one channel), where each value indicates how much a small change in that pixel would affect the model's output. High values represent pixels to which the model is very sensitive, so they are relevant for the prediction. Low values represent pixels that barely influence the model's decision. The resulting map can be viewed as a grayscale image (or overlaid with color) that highlights the areas of greatest influence.

Some cons are that sometimes it can be noisy and we do not have a sense of direction (positive or negative impact for the prediction). However, it is easy to implement and computationally efficient.

\subsection{Input x Gradient}

\textbf{Input x Gradient} is an attribution method that seeks to identify which parts of the input contribute most to the model's prediction by combining the magnitude of the input with the model's sensitivity to that input. The intuition is simple: if a pixel has a high value and the gradient (derivative of the output with respect to that pixel) is also high, then that pixel has a strong influence on the prediction.

The mathematical definition is very similar to the saliency map:
$$S_{i, j}(x)=\max_{k\in\{1, 2, 3\}}\left|x_{i, j, k}\cdot \dfrac{\partial f_c(x)}{\partial x_{i, j, k}}\right|.$$

It has the same pros and cons as the saliency map. In addition to this, sometimes it can be more informative than the saliency map, especially for inputs with many zeros or low intensity, but it always penalizes pixels with low values, and this can sometimes be incorrect.

\subsection{Guided Backpropagation}

\textbf{Guided backpropagation} is a gradient-based technique that aims to generate sharper and less noisy saliency maps than pure gradients. It modifies the backpropagation process to block negative gradients in the ReLU layers during activation and backpropagation. Thus, only positive gradients propagate through neurons that were also active in the forward direction. An example is shown in Figure \ref{backprop}.

The definition of $S(x)$ is the same as in the saliency maps, but with the modified backpropagation of ReLU layers.

This method is computationally efficient and usually produces cleaner, more focused saliency maps. However, it is not faithful to the true attribution: it does not respect the causal structure of the model.

\subsection{DeconvNet}

\textbf{DeconvNet} is a technique originally proposed to visualize which parts of an image trigger the activation of certain neurons in a CNN. Although its name suggests ``deconvolution'', it does not actually mathematically reverse convolutions, but rather performs modified backpropagation to produce interpretable visualizations. Unlike Guided Backpropagation, which requires both the forward activation and the backward gradient to be positive, DeconvNet only filters on the incoming gradient. This allows us to visualize which input patterns excite an internal neuron, regardless of whether the neuron was active in the forward pass. Backpropagation process is shown in Figure \ref{backprop}.

\begin{figure}[H]
\centering
\includegraphics[scale=0.3]{images/chapter_5/backprop}
\caption{Differences in ReLU backpropagation for saliency map, DeconvNet and guided bakpropagation.}
\label{backprop}
\end{figure}

Figure \ref{backprop} summarizes the differences in the backward pass of ReLUs using different methods:
\begin{itemize}
\item The first part referes to the forward pass of a ReLU layer. Negative values are set to zero.
\item The second image is the standard backpropagation. Gradient only propagates in positions that were not zero in the forward pass.
\item The third image is the backpropagation for DeconvNet. Gradient can propagate in positions that were zero in the forward pass, but now only positive gradients propagate.
\item The fourth image is the backpropagation for guided backpropagation. It is the same as standard backpropagation, but now only positive gradients propagate. It can be seen as the combination of standard backpropagation plus DeconvNet backpropagation.
\end{itemize}

\subsection{SmoothGrad}

\textbf{SmoothGrad} is an explainability technique that improves the quality of saliency maps generated by gradient methods, such as the classic saliency map. Its goal is to reduce noise and make visualizations more interpretable. It can be applied to any gradient-based method.

Direct gradient maps are often noisy and difficult to interpret due to the model's high sensitivity to small input perturbations. SmoothGrad reduces this noise by averaging gradients computed over multiple noisy versions of the original image.

\begin{figure}[H]
    \centering
    \begin{minipage}{0.43\textwidth}
        \centering
        \includegraphics[width=\linewidth]{images/chapter_5/perro}
    \end{minipage}
    \hfill
    \begin{minipage}{0.47\textwidth}
        \centering
        \includegraphics[width=\linewidth]{images/chapter_5/explanare}
    \end{minipage}
\caption{SmoothGrad applied to a dog. The original image is on the left and the explanation map on the right.}
\label{explanare}
\end{figure}

The \textbf{algorithm} is the following. Given a model $f$, an image $x$ and a class $c$:
\begin{enumerate}
\item $N$ perturbed images are generated with gaussian noise: $x_i = x + \mathcal{N}\left(0, \sigma^2\right) ~ \forall ~ i = 1, \ldots, N$.
\item For each perturbed image $x_i$, its explanation $S_i(x)$ is calculated with the desired method.
\item Return $S(x)=\dfrac{1}{N}\sum_{i=1}^NS_i(x).$
\end{enumerate}
Typically $N$ is between 10 and 100 and $\sigma$ is between the 10\% and 20\% of the range of pixel values (e.g., if the pixels are in $[0, 1]$, then $\sigma$ is between 0.1 and 0.2).

It reduces the noise of other techniques, but it is more computationally expensive, as $N$ forward and backward passes need to be done. An example is shown in Figure \ref{explanare}. 
 
\subsection{Integrated Gradients}

\textbf{Integrated Gradients (IG)} is a gradient-based explainability method that seeks to attribute the importance of each pixel (or input feature) to the model's prediction more robustly than direct gradients. As SmoothGrad, it can also be applied to any gradient-based method.

IG integrates gradients along a path from a base reference (e.g., a black image) to the original image. Let $f$ be a model, $x$ an input, and $x'$ the reference. The explanation is defined as:
$$\text{IG}(x)=(x-x')\int_{\alpha=0}^1 \dfrac{\partial f(x'+\alpha(x-x'))}{\partial x}d\alpha.$$
However, in practice we can not integrate, but remember from calculus that integrals are just sums! Therefore, we can redefine:
$$\text{IG}(x)=(x-x')\dfrac{1}{M}\sum_{k=1}^M\dfrac{\partial f\left(x'+\frac{k}{M}\left(x-x'\right)\right)}{\partial x},$$
where $M$ is the number of steps in the interpolation, typically between 50 and 300. The part of multipyling by $(x-x')$ is for the attribution to satisfy the completeness axiom, i.e., that the sum of all values will be equal to the difference between the original image and the baseline.

This is a very strong and robust method, but it is more computationally expensive. Moreover, the bigger $M$, the better it aproximates the integral, but again the more computationally expensive is.

The general idea is to average the gradients obtained from the baseline image to the original image and multiply this by the difference of both images. Figure \ref{ig} shows an example of an interpolation.

\begin{figure}[H]
\centering
\includegraphics[scale=0.53]{images/chapter_5/ig}
\caption{Interpolation of an image for different values of $\alpha$. The idea is to compute the saliency map for each interpolated image and then return the averaged results.}
\label{ig}
\end{figure}

\section{Ad-Hoc Methods}

Ad-hoc methods are techniques specifically designed to exploit the \textbf{internal structure} of certain models, especially deep convolutional networks (CNNs). Unlike general gradient or perturbation-based methods, these approaches rely on specific layers and operations of the network to generate interpretable visualizations, typically in the form of heat maps over the input image.

They are \textbf{model-specific} methods, but more restrictive than gradient-based methods, they can only be applied to a subset of CNNs that have some particular structure. They are also much more \textbf{efficient} than occlusion methods, as usually only one forward pass and one backward pass is needed.

\subsection{CAM}

\textbf{CAM} (Class Activation Mapping) is a visualization technique that allows you to identify \textbf{which regions of an image have been most relevant for a convolutional network to make a class prediction}. CAM requires an architecture with:
\begin{itemize}
\item A final convolutional layer followed directly by
\item A GAP (Global Average Pooling) layer and a linear classification layer.
\end{itemize}

Let's suppose a CNN that, before the classification layer, produces a set of $n$ activation maps $A_k(x)$ with shape $h\times w$, where $k=1, \ldots, n$ indexes the channels. Then the network applies:
\begin{enumerate}
\item Global Average Pooling:
$$F_k(x) = \dfrac{1}{hw}\sum_{i=1}^h\sum_{j=1}^wA_{i, j, k}(x) ~ \forall k=1, \ldots, n.$$
\item Classification layer (output for class $c$):
$$f_c(x) = \sum_{k=1}^n w^c_kF_k(x),$$
where $w_k^c$ is the learned weight that connects channel $k$ with class $c$.
\end{enumerate}
Doing some math, we have that
$$f_c(x) = \sum_{k=1}^n w^c_k\left(\dfrac{1}{hw}\sum_{i=1}^h\sum_{j=1}^wA_{i, j, k}(x)\right) = \dfrac{1}{hw}\sum_{i=1}^h\sum_{j=1}^w\sum_{k=1}^n w_k^cA_{i, j, k}(x).$$
This suggests that a \textbf{class activation map} $M_c(i, j)$ can be constructed as:
$$M_c(i, j) = \sum_{k=1}^n w_k^cA_{i, j, k}(x).$$

This map indicates the importance of each spatial position $(i, j)$ for class $c$. By \textbf{interpolating} it back to the original image size, we obtain an interpretable visualization of ``where the network looks'' to classify an image as class $c$.

This method is intuitive and efficient, but the most obvious limitation is that it can only be applied to a few networks. Finally, a graphical representation is shown in Figure \ref{cam}.

\begin{figure}[H]
\centering
\includegraphics[scale=0.4]{images/chapter_5/cam}
\caption{Graphical representation of CAM.}
\label{cam}
\end{figure}


\subsection{Grad-CAM}

\textbf{Grad-CAM} is a \textbf{generalization of CAM} that allows activation maps to be generated in any modern CNN, without requiring a special architecture. Its goal remains the same: to visualize which image regions have most influenced the prediction of a given class.

Given a model $f(x)$, an input image $x$, and a class of interest $c$, proceed as follows:
\begin{itemize}
\item Select a convolutional layer (usually the last one) and obtain its activation maps $A_k(x)$, where $k$ is the channel (feature map). We will assume there are $n$ channels and the shape of the activation is $h\times w$.
\item Calculate the gradients of the model output for class $c$ with respect to these activation maps:
$$\dfrac{\partial f_c(x)}{\partial A_k(x)} ~ \forall k = 1, \ldots, n.$$
\item Average those gradients over the $(i, j)$ spatial space to obtain a weight per channel:
$$\alpha_k^c = \dfrac{1}{hw}\sum_{i=1}^h\sum_{j=1}^w\dfrac{\partial f_c(x)}{\partial A_{i, j, k}(x)}.$$
\item Combine the activated maps using these weights (ReLU ensures that only regions positively associated with the $c$ class are displayed):
$$M_c(i, j) = \text{ReLU}\left(\sum_{k=1}^n\alpha_k^cA_{i, j, k}(x)\right).$$
\end{itemize}
$M_c$ is a heat map of the same resolution as the activation maps. By interpolating it to the original image size, it can be overlaid to visualize relevant regions.

This method is also intuitive (it is like a variant of CAM) and computationally efficient, and it can be applied to any convolutional network. However, it can not be applied to traditional neural networks, for example when we work with features, not images.

\subsection{Guided Grad-CAM}

\textbf{Guided Grad-CAM} is a technique that merges two popular explainability methods in convolutional networks:
\begin{itemize}
\item Grad-CAM, which generates a localized, but blurry (low-detail) activation map.
\item Guided backpropagation, which produces a fine-grained, high-resolution map with little contextual information.
\end{itemize}
The idea is to multiply both maps to obtain a visualization that is spatially accurate, detailed at the pixel level and relevant to the predicted class.

Guided Grad-CAM is simply the element-by-element multiplication of both the maps obtained by Grad-CAM and guided backpropagation. Grad-CAM tells you which regions are important, guided backprop tells you which pixels are important and, therefore, guided Grad-CAM tells you which important pixels are in the relevant regions.

\subsection{Attention as Explainability}

Some models, like Vision Transformers (ViTs), use \textbf{attention} on the inputs to obtain the final prediction. Although it may be tempting to use attention weights as importance maps, there are several reasons why this is problematic:
\begin{itemize}
\item They are not necessarily causal: high attention weights do not imply that one part of the input causes the output. They can be spurious correlations or training artifacts.
\item They are context-dependent: in models with multiple heads and layers, it is difficult to attribute meaning to a single attention map without considering more complex interactions.
\end{itemize}

Attention can be useful in some simple models, such as visual attention in classification or captioning tasks. In more complex models, if accompanied by sensitivity analyses, counterfactuals, or intervention techniques, attention maps can be a clue.

In conclusion, attention is not explanation and \textbf{``saying that attention is explainability is like saying that correlation is causation''}. It's an accurate analogy: both can be clues, but they do not guarantee true interpretability. Using attention as an explanation requires caution, additional validation, and a clear understanding of its role in architecture. 

\begin{figure}[H]
\centering
\includegraphics[scale=0.45]{images/chapter_5/techniques}
\caption{Different explainability techniques.}
\label{techniques}
\end{figure}

To end this section Figure \ref{techniques} shows some techniques for three inputs. Note that the maps are shown themselves alone instead of displayed above the original image.

\section{Evaluation}

It is not enough to generate eye-catching visualizations: an explainability technique must be reliable, consistent and useful. To achieve this, it is essential to evaluate its quality using objective metrics or specific tests. This evaluation can be approached from several angles This section presents the main tools and criteria used to assess the quality of an explanation, both quantitatively and qualitatively.

\subsection{Common Problems}

\begin{itemize}
\item \textbf{Unstable gradients}: Gradients can be extremely sensitive to small changes in the input. This means that small movements in a pixel value can lead to large changes in the magnitude or even flip the sign of a graadient. This can lead to saliency maps that are hard to interpret, as they appear speckled or random in regions that should probably be smooth or semantically meanningful. SmoothGrad has been shown to help. It usually reveals more consistent patterns.

\item \textbf{Rapid sign swings}: A related problem is that we can often see rapid swings in the signs of the gradients. This has to do both with the sensitivity of the gradients and the nonlinearity of neural networks. Ultimately, this emphasizes the points that we should be taking the absolute value of the gradients. This is especially true when you want to aggregate gradients over different regions of the image.

\item \textbf{Irrelevant information}: The previous two problems could be summarized as one, and that is noisy gradients. Another factor that contributes to noisy saliency maps is that often gradients can contain information that is not specific to the target logit. This can come from pixels in the background or those from features associated with other classes. Guided backpropagation can clean up saliency maps using the ReLU masking, as it suppresses negative gradients during backpropagation, leading to a visually sharper saliency map that highlithts the features used by a model. However, this method is not class discriminative, meaning it may still highlight features useful for multiple classes, not just the target one. Grad-CAM is a more reliable class discriminative method. We can combine both methods with Guided Grad-CAM.

\item \textbf{Saturated gradients}: In some cases, the gradients of important pixels can still be small or even zero. During backpropagation, if the input to the ReLU function has a negative value, then the gradients will be zero, regardless of how large the negative input is. Integrated Gradients can solve this problem, because we observe how gradients change over a path.

\item \textbf{Vanishing gradients}: In large networks, the gradients of earlier layers can become small because they are obtained by iteratively multiplying gradients of deepeer layer, which are small values themselves. Methods like Grad-CAM, that only use gradients of deeper layers, can solve this problem.

\item \textbf{Lack of context}: Gradients reflect local sensitivity. That is, how a small change in an input pixel would impact the output. This does not capture how specific regions of the image contribute to a prediction. Integrated Gradients and SmoothGrad can solve this, because it changes all pixels at the same time. Grad-CAM also mitigates this problem, as it usually relies on deeper layers, which contains simpler features and more informative pixels.
\end{itemize}

\subsection{Loyalty and Inverse Loyalty}

The loyalty and inverse loyalty metrics are two popular ways to evaluate the fidelity of an explainability technique. That is, how well it truly reflects the importance of the identified regions relative to the model's prediction.

\textbf{Loyalty} measures how well the regions identified as important by an explanation match what actually influences the model's prediction.

The intuition is that if you remove what the salience map says is important, the prediction should change significantly. A good explanation will cause a rapid drop in the predicted probability for the target class. This is usually represented as a curve: percentage of removed pixels vs. difference with the model output. To quantify it, the area under the curve (AUC) is typically used. The closer to 1, the better, because it means that the most important pixels identified were actually important for the prediction.

The procedure is the following:
\begin{enumerate}
\item Start with the original image $x$ and its prediction $f_c(x)$ for class $c$.

\item The pixels (or regions) are ordered from highest to lowest importance according to the explanation map.

\item These regions are progressively removed (e.g., by setting them to zero).

\item For each generated image, we record how the model's output changes for class $c$. For example, we can define it as
$$1-\left|\dfrac{f_c(x)-f_c(x')}{f_c(x)}\right|,$$
where $f_c(x)$ is the original prediction and $f_c(x')$ the prediction of the modified image.
\end{enumerate}

\textbf{Inverse loyalty} measures how well the regions identified as irrelevant by an explanation match what actually does not influence the model's prediction.

The intuition is that if you remove what the salience map says is less important, the prediction should not change. A good explanation will cause a slow drop in the predicted probability for the target class. This is usually represented as a curve. The closer the AUC to 0, the better, because it means that the less important pixels identified were actually irrelevant for the prediction. At the end, we should see a rapid drop in the probability, as important pixels will be removed.

The procedure is the following:
\begin{enumerate}
\item Start with the original image $x$ and its prediction $f_c(x)$ for class $c$.

\item The pixels (or regions) are ordered from lowest to highest importance according to the explanation map.

\item These regions are progressively removed (e.g., by setting them to zero).

\item For each generated image, we record how the model's output changes for class $c$. We can define the difference between predictions as in loyalty.
\end{enumerate}

\begin{figure}[H]
    \centering
    \begin{minipage}{0.19\textwidth}
        \includegraphics[width=\linewidth]{images/chapter_5/metrics/a1}
    \end{minipage}
    \begin{minipage}{0.19\textwidth}
        \includegraphics[width=\linewidth]{images/chapter_5/metrics/a2}
    \end{minipage}
    \begin{minipage}{0.19\textwidth}
        \includegraphics[width=\linewidth]{images/chapter_5/metrics/a3}
    \end{minipage}
    \begin{minipage}{0.19\textwidth}
        \includegraphics[width=\linewidth]{images/chapter_5/metrics/a4}
    \end{minipage}
    \begin{minipage}{0.19\textwidth}
        \includegraphics[width=\linewidth]{images/chapter_5/metrics/a5}
    \end{minipage}
    \caption{Removing most important pixels of an image.}
\label{questions1}
\end{figure}

\begin{figure}[H]
    \centering
    \begin{minipage}{0.19\textwidth}
        \includegraphics[width=\linewidth]{images/chapter_5/metrics/b1}
    \end{minipage}
    \begin{minipage}{0.19\textwidth}
        \includegraphics[width=\linewidth]{images/chapter_5/metrics/b2}
    \end{minipage}
    \begin{minipage}{0.19\textwidth}
        \includegraphics[width=\linewidth]{images/chapter_5/metrics/b3}
    \end{minipage}
    \begin{minipage}{0.19\textwidth}
        \includegraphics[width=\linewidth]{images/chapter_5/metrics/b4}
    \end{minipage}
    \begin{minipage}{0.19\textwidth}
        \includegraphics[width=\linewidth]{images/chapter_5/metrics/b5}
    \end{minipage}
    \caption{Removing less important pixels of an image.}
\label{questions2}
\end{figure}

\begin{figure}[H]
    \centering
    \begin{minipage}{0.49\textwidth}
        \includegraphics[width=\linewidth]{images/chapter_5/metrics/af}
    \end{minipage}
    \begin{minipage}{0.49\textwidth}
        \includegraphics[width=\linewidth]{images/chapter_5/metrics/bf}
    \end{minipage}
    \caption{Loyalty (left) and inverse loyalty (right) curves.}
\end{figure}

\subsection{ROAR}

The \textbf{ROAR (RemOve And Retrain)} metric is a technique for evaluating the explainability of machine learning models, especially deep neural networks. Its objective is to measure how useful the generated explanations (e.g., heatmaps of feature or pixel importance) are in understanding the model's decisions. ROAR does not directly measure the explanation, but rather its usefulness: if removing what is ``important'' according to the explanation causes the model to suffer, then the explanation was useful in detecting what was relevant. The biggest disadvantage is that it is computationally expensive.


These steps are followed:
\begin{enumerate}
\item You train an original model with a dataset.

\item You generate the explainability maps with a technique.

\item For each image, you remove the most important pixels according to each explanation. For example, you nullify the most important pixels (setting them to black, zero, or some neutral value).

\item You train a new model from scratch with this ``corrupted'' data.

\item If the explanation technique is good, then the new model should perform worse because it lost crucial information. A drop in accuracy is interpreted as a sign of good explainability.
\end{enumerate}

Figure \ref{roar} shows an example on what is not a good explainability technique. Orange represents the pixels removed in the order according to the explainability technique, black is a random order to use as baseline and blue the inverse of orange (from less important to most important pixels):
\begin{itemize}
\item If we only look at the left curve, we may think the technique is good. As we remove pixels accuracy decreases. Moreover, removing the most important pixels (orange) causes the greatest degradation in the accuracy, the less important pixels (blue) the least degradation, and removing pixels randomly (black) is an intermediate point.

\item However, looking at the right curve, if we remove the same pixels, but retrain the model in each case, the accuracy is greater at first removing the most important pixels (according to the explainability technique) than random pixels! That makes as think that the technique may be is not accurate.
\end{itemize}

In the original paper, no quantitative measures are given. One could be the relative worsening in accuracy, f1 or the desired metric and finally take the area under the curve.

\begin{figure}[H]
\centering
\includegraphics[scale=0.65]{images/chapter_5/roar}
\caption{Using ROAR, we can see that after re-training most of the information is actually still present.}
\label{roar}
\end{figure}

\subsection{Sanity Checks}

\textbf{Sanity checks} for saliency maps are a set of tests designed to verify whether the explanations generated by interpretability methods truly reflect what the model has learned, or whether they are actually independent of the trained model and therefore useless. Many methods produce beautiful maps that appear to highlight important regions of the image, but there is no guarantee that they actually rely on the model's learned weights. Some methods can give similar results even if the model is untrained.

There are two main checks:
\begin{itemize}
\item \textbf{Model parameter randomization test}: The layers of the trained model are progressively replaced with random weights (from output to input). The saliency map should change substantially as the model's knowledge is destroyed. If the map does not change much even though the model is randomized, it means the explanation does not depend on the trained model, so it is a bad explanation!
\item \textbf{Data randomization test}: A model is trained with random labels and saliency maps are then generated. If the explanation method produces maps similar to those of the well-trained model, then it is not capturing relevant information, so it is unreliable.
\end{itemize}

Figure \ref{check1} shows an example of the model parameter randomization test. Some methods produce the same output when randomizing almost all layers. That is an indicator that the technique may be is not correct.

Figure \ref{check2} shows the data randomization test on a CNN trained on the MNIST dataset. In the left part, we can see examples of explanations for different methods to a image with digit zero. In some cases the explanation is almost the same. Again, that is an indicator that the technique may be is not correct. To have a more quantitative number, in the right part is shown the correlation between the explanations when training with true and random labels. In this particular case, Grad-CAM seems to be the best technique, while Integrated Gradients and Gradient x Input seem to be the worst.

\begin{figure}[H]
\centering
\includegraphics[scale=0.87]{images/chapter_5/check2}
\caption{Model parameter randomization test.}
\label{check1}
\end{figure}

\begin{figure}[H]
    \centering
    \begin{minipage}{0.55\textwidth}
        \includegraphics[width=\linewidth]{images/chapter_5/check1a}
    \end{minipage}
    \begin{minipage}{0.44\textwidth}
        \includegraphics[width=\linewidth]{images/chapter_5/check1b}
    \end{minipage}
\caption{Data randomization test.}
\label{check2}
\end{figure}

\section{Exercises}

\question{}{
Give some arguments about why explainability is important. Think about particular situations where it can be useful.
}

\question{}{
Write a program to load a pretrained network such as ResNet and donwload some images, like cats and dogs. We will use the network and the custom dataset for the following exercises.
}

\question{}{
Write a program to obtain the activations from different layers in a CNN. Plot them to see how they evolve as they approach deeper layers.
}

\question{}{
Write a program to see the maximally activating patches of an image for all its layers.
}

\question{}{
Generate an image that maximizes the probability for ``dog'' class using gradient ascent.
}

\question{}{
Generate an image using feature inversion for the layer you want in a CNN. The original image can belong, for example, to the ``cat'' class.
}

\question{}{
Use DeepDream in the image you want in different layers. What can you see?
}

\question{}{
Write the code for the approximation of Shapley values (the common approach followed in the practice). Generate a map for the most important pixels in the image you want.
}

\question{}{
Write the code for the occlusion method. Generate a map for the most important pixels in the image you want. If the network is heavy, it may take some minutes. Take into account that bigger patch size implies less computation time.
}

\question{}{
Write the code for RISE. Generate a map for the most important pixels in the image you want. If the network is heavy, it may take some minutes. Also, take into account that bigger patch size implies less computation time.
}

\question{}{
Write the code for saliency map. Generate a map for the most important pixels in the image you want.
}

\question{}{
Write the code for Input x Gradient. Generate a map for the most important pixels in the image you want.
}

\question{}{
Write the code for guided backpropagation. Generate a map for the most important pixels in the image you want. As you need to modify the backward pass, you will need to use PyTorch hooks. If you don't remember them well, I encourage you to see these two videos:
\begin{itemize}
    \item \href{https://www.youtube.com/watch?v=syLFCVYua6Q}{PyTorch Hooks Explained - In-depth Tutorial}
    \item \href{https://www.youtube.com/watch?v=1ZbLA7ofasY}{Visualizing activations with forward hooks (PyTorch)}
\end{itemize}
}

\question{}{
Write the code for DeconvNet. Generate a map for the most important pixels in the image you want. You will also need to make use of hooks.
}

\question{}{
Write the code for SmoothGrad. Choose any method and generate a map for the most important pixels in the image you want.
}

\question{}{
Write the code for integrated gradients. Generate a map for the most important pixels in the image you want.
}

\question{}{
Write the code for CAM. Generate a map for the most important pixels in the image you want. You will need to make use of hooks.
}

\question{}{
Write the code for Grad-CAM. Generate a map for the most important pixels in the image you want. You will need to make use of hooks.
}

\question{}{
Write the code for Guided Grad-CAM. Generate a map for the most important pixels in the image you want. Note that if you have already coded guided backpropagation and Grad-CAM it is very straightforward.
}

\question{}{
Write the code to compute loyalty and inverse loyalty and its AUC. Optionally, write the code to see the evolution of the image when you remove the top or less important pixels according to the explanation, like in Figures \ref{questions1} and \ref{questions2}.
}

\question{}{
Write the code to compute the ROAR curve. You can use a simple dataset, as it may take some time for retraining.
}

\question{}{
Write the code to generate the model parameter randomization test and the data randomization test. Also, try to replicate Figures \ref{check1} and \ref{check2}.
}
